\documentclass[11pt]{article}

\usepackage[margin=1in]{geometry}
\usepackage[T1]{fontenc}
\usepackage[utf8]{inputenc}
\usepackage{lmodern}
\usepackage{microtype}
\usepackage{amsmath, amssymb, amsthm}
\usepackage{booktabs}
\usepackage{array}
\usepackage{graphicx}
\usepackage{xcolor}
\usepackage{hyperref}
\usepackage{enumitem}
\usepackage{caption}
\usepackage{subcaption}
\usepackage{longtable}
\usepackage{float}

\hypersetup{
    colorlinks=true,
    linkcolor=blue!50!black,
    citecolor=blue!50!black,
    urlcolor=blue!50!black
}

\setlist[itemize]{leftmargin=1.5em}
\setlist[enumerate]{leftmargin=1.5em}

\newcommand{\hphase}{h_1(T)}
\newcommand{\modsector}{s_{6,M}}
\newcommand{\zprimary}{2.575}
\newcommand{\zsensitivity}{2.326}
\newcommand{\Arc}{\operatorname{Arc80}}
\newcommand{\R}{\mathcal{R}}
\newcommand{\Ssector}{S_{\mathrm{sector}}}
\newcommand{\Sint}{S_{\mathrm{int}}}
\newcommand{\Smain}{S_{\mathrm{main}}}
\newcommand{\principaldoi}{10.5281/zenodo.19827222}

\newcommand{\optionalfigure}[3]{%
\begin{figure}[H]
    \centering
    \IfFileExists{#1}{%
        \includegraphics[width=0.98\textwidth]{#1}%
    }{%
        \fbox{\parbox[c][2.25in][c]{0.88\textwidth}{\centering Figure placeholder: \texttt{#1}}}%
    }
    \caption{#2}
    \label{#3}
\end{figure}
}

\title{Orthogonal Row-Sector Rebound Dynamics in Terminal Zeta-Zero Spacing Residuals}
\author{Jeffery Huckstead\\Independent Researcher}
\date{Preprint draft; DOI to be assigned}

\begin{document}

\maketitle

\begin{abstract}
We identify a row-sector rebound structure in terminal adjacent-zero spacing residuals near a $\sim 10^3$-row sector boundary identified through an unbiased window-scale sweep, without prior specification of the active modulus. Under a base-6 row-sector decomposition, the structure is amplitude-dominant for modulus $1009$, while modulus $990$ remains co-significant under permutation guardrails. Follow-on audits reject simple $h_1$-corridor gating and do not confirm local $h_1 \times$ sector modulation. A synthetic null calibration shows that apparent growth of interaction strength under finer local-$h_1$ binning matches bin-resolution sparsity inflation rather than hidden phase-conditioned structure. Thus the row-sector rebound and the previously published $h_1$ packet organization are empirically distinct on the observed terminal horizon.
\end{abstract}

\noindent\textbf{Keywords:} Riemann zeta zeros; zero spacing; row-sector rebound; tail-sign polarity; empirical audit; terminal band; observed-horizon result.

\section{Introduction and Relation to Prior Work}

The principal paper established that rare extreme compressions in adjacent Riemann zeta-zero spacing organize in the first logarithmic harmonic
\begin{equation}
    h_1(T) = \left\{\frac{\log(T/2\pi)}{\pi}\right\},
    \label{eq:h1}
\end{equation}
stably across resolution ablation and with unusual tightening in the terminal observational band \cite{huckstead_h1_principal}. That result concerns the \emph{phase location} of compression packets: where in logarithmic phase space rare compressions prefer to land. The broader statistical context includes the classical zero-counting and spacing background of the Riemann zeta function \cite{riemann1859,montgomery1973,mehta_random_matrices}, and the computations build on tabulated zero ordinates such as those made available by Odlyzko \cite{odlyzko_tables}.

The present paper addresses a different question. Within the terminal residual tail population, how are compression and expansion events distributed across row-index sectors? This question is structurally distinct from the $h_1$ concentration question and requires a separate instrument. Rather than asking where compression packets cluster in a continuous logarithmic coordinate, we ask how terminal tail polarity is distributed across a discrete row-index decomposition.

The main result is a row-sector rebound structure: compression and expansion tails are not distributed identically across row-index sectors. Among the tested decompositions, modulus $1009$ is the most consistent sector-amplitude ruler, while modulus $990$ remains co-significant under permutation guardrails. Five follow-on audits establish that this sector structure is empirically orthogonal to the previously published $h_1$ organization on the observed terminal horizon: it is not gated by $h_1$ corridor membership, not confirmed as a local $h_1 \times$ sector interaction, and not explained by bin-resolution artifacts.

Table~\ref{tab:distinction} summarizes the relationship between the published work and the present continuation.

\begin{table}[H]
\centering
\caption{Structural distinction between the published paper and the present continuation.}
\label{tab:distinction}
\renewcommand{\arraystretch}{1.25}
\begin{tabular}{p{0.18\textwidth}p{0.38\textwidth}p{0.36\textwidth}}
\toprule
 & \textbf{Published paper} & \textbf{Present paper} \\
\midrule
Question & Where do rare extreme compressions cluster? & How is terminal tail polarity distributed across row-index sectors? \\
Coordinate & $h_1(T)$ logarithmic phase & Row-index sector modulo $M$, base-6 decomposition \\
Primary finding & $h_1$ phase concentration and narrow $\Arc$ corridor & Sector amplitude asymmetry, strongest under modulus $1009$ \\
Terminal effect & Compression packet tightening in the terminal band & Permutation-significant sector polarity and rebound structure \\
Relationship & Not applicable & Empirically orthogonal to $h_1$ organization under tested coupling forms \\
Horizon & Published $h_1$ corridor horizon & Terminal board used in P74--P78, with explicit observed-horizon limitation \\
\bottomrule
\end{tabular}
\end{table}

This paper does not claim that $h_1$ controls or gates the row-sector rebound, that local $h_1 \times$ sector modulation is confirmed, or that modulus $1009$ is uniquely proven as a physical modulus. The scope of each positive finding is bounded by the explicit null testing reported in Sections~\ref{sec:audit_chain} and~\ref{sec:results_summary}. The negative findings are part of the result: they prevent the continuation structure from being folded back into the $h_1$ packet-corridor claim.

\section{Data and Instrument}

\subsection{Source Board}

All computations use the seam-stitched terminal checkpoint parquet from the principal paper,
\begin{center}
\texttt{df\_master\_terminal\_checkpoint\_v16c\_8000\_31p5B.parquet},
\end{center}
containing $935{,}111$ rows with altitude range approximately $28{,}499.5\mathrm{M}$ to $30{,}610.0\mathrm{M}$. Throughout this paper, all claims are empirical claims on this observed terminal horizon unless explicitly stated otherwise.

\subsection{Residual Construction}

For a window of size $W$ centered at row midpoint $r_{\mathrm{mid}}$, the local spacing metric is computed from one of two residual-ratio coordinates:
\begin{align}
    \mathrm{Gap\_to\_Pure\_Ratio} &= \frac{\Delta_k}{s_{\mathrm{pure}}(T_k)}, \\
    \mathrm{Gap\_to\_Fitted\_Ratio} &= \frac{\Delta_k}{\widehat{s}(T_k)},
\end{align}
where $s_{\mathrm{pure}}$ is the pure logarithmic spacing spine and $\widehat{s}$ is the fitted spine used in the principal audit chain. For each window family, altitude drift is removed by selecting a best-BIC baseline from the audit model family, and residuals are converted to robust $z$-scores.

The primary tail threshold is
\begin{equation}
    |z| \geq \zprimary,
    \label{eq:primary_tail}
\end{equation}
with a sensitivity threshold
\begin{equation}
    |z| \geq \zsensitivity.
    \label{eq:sensitivity_tail}
\end{equation}
Compression tails satisfy $z \leq -\zprimary$ and expansion tails satisfy $z \geq +\zprimary$ in the primary analysis.

\subsection{Row-Sector Coordinate}

For modulus $M$ and window row midpoint $r_{\mathrm{mid}}$, define the base-6 row sector
\begin{equation}
    s_{6,M}(r_{\mathrm{mid}})
    = \left\lfloor 6 \cdot \frac{r_{\mathrm{mid}} \bmod M}{M} \right\rfloor
    \in \{0,1,2,3,4,5\}.
    \label{eq:sector}
\end{equation}
The tested moduli are
\begin{equation}
    M \in \{990, 997, 1009, 1080\}.
\end{equation}
The primary stride is \texttt{stride\_mod6}; the sensitivity stride is \texttt{stride\_W\_4}. The window range used for the final calibration chain is
\begin{equation}
W \in \{900,925,930,950,975,1000,1009,1050,1100,1125,1200,1250,1300\}.
\end{equation}

\subsection{Frozen Parameter Registry}

The frozen registry records the altitude split point $30{,}000\mathrm{M}$, primary and sensitivity tail thresholds, moduli tested, stride modes, local $h_1$ bin counts, P78 cell-count floors, and the literal-claim matrix rule. In particular, P78 uses a minimum mean cell-count floor of $15$ and a minimum cell count of $5$ when interpreting interaction ratios. These parameters are recorded in the patched frozen parameter registry and P79B manifest (timestamp \texttt{\p79btimestamp}) before final paper-language assignment.

\section{Audit Chain}
\label{sec:audit_chain}

The empirical case is built through five sequential audits, P74--P78, followed by the P79/P79B claim freeze. Each audit uses a pre-specified read rule and a separate null or calibration structure. Figure~\ref{fig:literal_matrix} presents the final literal-claim evidence matrix. Each cell answers whether the audit supports the claim as written, rather than whether the audit supports the final ledger disposition.

\optionalfigure{p79b_literal_claim_evidence_matrix.png}{P79B literal claim evidence matrix. Each cell answers: does this audit support the literal claim as written? For rejected claims, an audit that contradicts the claim is coded as \emph{rejects}, not as support for the rejection.}{fig:literal_matrix}

\subsection{P74: Unbiased Window-Scale Rebound Sweep}

P74 conducted an unbiased sweep over window sizes and sector configurations to identify whether a rebound seam exists without pre-specifying the active modulus. The moduli $990$, $997$, $1009$, and $1080$ were identified as the strongest candidates from the P74 sweep result and then carried forward without reselection. The strongest permutation-tested cluster occurred near the $\sim 10^3$-row scale, with top configurations concentrated in the $W \approx 925$--$1300$ band under base-6 sectorization. The audit verdict was
\begin{center}
\texttt{BALANCE\_SEAM\_REBOUND\_PEAK\_SURVIVES\_SWEEP}.
\end{center}
This established the rebound seam as a candidate structure, but did not by itself determine whether the structure shares a chassis with the previously published $h_1$ organization.

\subsection{P75: $h_1$-Gated Rebound Audit}

P75 tested whether the row-sector rebound is stronger inside the published $h_1$ corridor than in shoulder regions. Let
\begin{equation}
    \Delta_{h_1} = \mathrm{rebound\_score}_{\mathrm{core}} - \mathrm{rebound\_score}_{\mathrm{shoulder}}.
\end{equation}
The corridor was loaded from the published manifest without retuning. Four nulls were used: sector shuffle within $h_1$ bins, $h_1$-bin shuffle within sector and altitude bins, circular row-phase shift, and altitude-matched $h_1$ assignment. Across the tested configurations, $\Delta_{h_1}$ did not survive the null battery. The audit verdict was
\begin{center}
\texttt{ORTHOGONAL\_ROW\_REBOUND\_CONFIRMED}.
\end{center}
Thus, the row-sector rebound is not supported as a simple function of membership in the published $h_1$ corridor.

\subsection{P76: Global $h_1 \times$ Sector Decomposition}

P76 constructed a two-dimensional imbalance map over global $h_1$ bins and base-6 row sectors. For $h$-bin $a$ and sector $s$, define the tail-sign imbalance
\begin{equation}
    I(a,s) = \log\left(\frac{E(a,s)+1/2}{C(a,s)+1/2}\right),
    \label{eq:imbalance}
\end{equation}
where $E(a,s)$ and $C(a,s)$ are expansion and compression counts in the cell. The interaction residual is
\begin{equation}
    R(a,s) = I(a,s) - \overline{I}(a,\cdot) - \overline{I}(\cdot,s) + \overline{I}.
    \label{eq:interaction_residual}
\end{equation}
Here $\overline{I}(a,\cdot)$ denotes the mean of $I(a,s)$ over $s$, $\overline{I}(\cdot,s)$ denotes the mean over $a$, and $\overline{I}$ denotes the grand mean. The strength ratio $\Sint/\Smain$ was used to compare interaction residual magnitude to the main-effect scale.

The important diagnostic from P76 was structural: the terminal tail population was so tightly concentrated in global $h_1$ that most global $h_1$ bins were empty. Consequently, the global interaction map had no vertical leverage, and $\Sint/\Smain$ collapsed to numerical zero. The audit verdict was
\begin{center}
\texttt{SECTOR\_DOMINANT\_H1\_WEAK\_MODULATOR}.
\end{center}
P76 therefore motivated a local-$h_1$ instrument rather than a global-bin instrument.

\subsection{P77: Local-$h_1$ Leverage and Sector Engine Audit}

P77 replaced global $h_1$ bins with local quantile bins computed within the terminal tail support. The primary local instrument used four quantile bins; the sensitivity instrument used six. This restored vertical leverage while preserving the sector coordinate from Equation~\eqref{eq:sector}.

The sector engine survived tail-label shuffle for the leading moduli, especially moduli $1009$ and $990$. The local $h_1 \times$ sector interaction residual did not survive null testing. The mod1009 consistency check ranked modulus $1009$ first by $\Ssector$ amplitude in the tested primary lanes. The audit verdict was
\begin{center}
\texttt{MOD1009\_SECTOR\_ENGINE\_CONFIRMED\_H1\_ORTHOGONAL}.
\end{center}
The P77 local imbalance maps are useful descriptive objects, but they are not promoted to confirmed local-$h_1$ interaction claims because their interaction component did not survive the null battery.

\subsection{P78: Bin-Resolution Null Calibration}

P78 tested the apparent rise in $\Sint/\Ssector$ under finer local-$h_1$ binning. The bin count was swept over
\begin{equation}
    B \in \{2,3,4,5,6,8,10,12\},
\end{equation}
with matched tail-label nulls and a synthetic Bernoulli-null control on the real geometry. The key observation was that the synthetic null produces the same rising interaction-ratio shape as bin count increases. Thus the 4-to-6-bin rise is a sparsity and bin-resolution effect, not a hidden local-$h_1$ signal.

The P78 primary calibration curve is shown in Figure~\ref{fig:p78_curve}. Observed curves remain within the null envelope, and polynomial tests do not show significant excess intercept or curvature for the primary mod1009 lane. The audit verdict was
\begin{center}
\texttt{BIN\_RESOLUTION\_INFLATION\_CONFIRMED},
\end{center}
with the secondary finding
\begin{center}
\texttt{MOD1009\_AMPLITUDE\_LEADER\_CONFIRMED}.
\end{center}

\optionalfigure{p78_primary_mod1009_calibration_curve.png}{P78 primary calibration curve for the mod1009, \texttt{stride\_mod6} lane. The observed local-$h_1$ interaction-ratio curve remains within the matched null envelope, indicating bin-resolution inflation rather than confirmed local-$h_1$ modulation.}{fig:p78_curve}

\section{Results Summary}
\label{sec:results_summary}

The final claim hierarchy is generated by the P79B hardened ledger. It separates publishable claims, supplemental claims, and rejected or quarantined claims.

\subsection{Level 1: Confirmed Claims}

\paragraph{C01. Row-sector rebound structure exists.}
Terminal adjacent-zero spacing residuals exhibit a row-sector rebound structure: compression and expansion tails are not distributed identically across row-index sectors. This claim is supported by the unbiased rebound sweep in P74, the sector-engine null tests in P77, and the P78 amplitude guardrail.

\paragraph{C03. Modulus 1009 is the most consistent sector-amplitude ruler.}
Among the tested row-sector decompositions, modulus $1009$ is the most consistent sector-amplitude ruler across $32$ of $32$ tested configurations spanning two metrics, two stride modes, and eight local bin counts, while modulus $990$ remains co-significant by permutation guardrail. This does not prove modulus $1009$ is the unique physical modulus; p-value rankings are ceiling-bounded and modulus $990$ remains co-significant.

\paragraph{C09. The $h_1$ packet corridor and row-sector rebound are empirically distinct structures.}
The previously published $h_1$ packet organization and the row-sector rebound are empirically distinct on the observed terminal horizon. The rebound is not gated by corridor membership, not confirmed as a local $h_1 \times$ sector modulation, and not explained by bin-resolution artifacts. Distinct on the observed terminal horizon does not preclude a deeper common mechanism at greater heights or under a different mathematical framework. The current evidence rules out the tested coupling forms; it does not rule out untested ones.

\subsection{Level 2: Supported Claims}

\paragraph{C02. The rebound is active near a roughly $10^3$-row seam.}
The rebound structure is most naturally studied near a roughly $10^3$-row sector boundary identified through an unbiased window-scale sweep, without prior specification of the active modulus. This is an observed-horizon scale, not a universal correlation length theorem.

\paragraph{C04. Modulus 990 remains co-significant but lower-amplitude than modulus 1009.}
Modulus $990$ remains co-significant under permutation guardrails, but modulus $1009$ carries larger sector amplitude on the observed horizon. P-value ceiling effects prevent unique ranking by significance alone.

\subsection{Level 4: Rejected or Quarantined Claims}

The following claims were explicitly tested and rejected, or quarantined as too strong for the available evidence:
\begin{itemize}
    \item \textbf{C05.} Global $h_1$ corridor membership gates row-sector rebound. This is rejected by P75 and contradicted by subsequent interaction audits.
    \item \textbf{C06.} Local $h_1 \times$ sector modulation is confirmed. This is rejected or unconfirmed by P76--P78.
    \item \textbf{C07.} The 4-to-6-bin rise is hidden $h_1$ structure. This is rejected by the P78 synthetic and matched null calibration.
    \item \textbf{C08.} Modulus $1009$ is uniquely proven as the physical modulus. This is quarantined. Amplitude leadership is confirmed; uniqueness is not.
\end{itemize}

\section{Discussion}

The audit chain resolves a tension raised by the principal paper. Does the terminal band's strong $h_1$ phase concentration reflect a single organizing mechanism that also governs compression/expansion polarity, or are there two distinct structures sharing the same terminal real estate?

The evidence favors the second interpretation. The $h_1$ organization and the row-sector rebound coexist in the terminal band, but they do not appear to share the tested coupling forms. The $h_1$ coordinate organizes packet location in a continuous logarithmic phase. The row-sector instrument detects a discrete tail-sign polarity structure across row-index sectors. These are orthogonal empirical structures on the observed terminal horizon.

This negative result is structurally informative. It prevents the sector rebound from being misdescribed as merely a shadow of the $h_1$ corridor. It also prevents the amplitude leadership of modulus $1009$ from being promoted into a uniqueness claim. The resulting model is narrower but stronger: a terminal row-sector rebound exists; modulus $1009$ is the most consistent amplitude ruler; modulus $990$ remains co-significant; and the tested $h_1$ coupling forms fail.

\subsection{What Would Change These Conclusions}

Disambiguation of modulus $1009$ and modulus $990$ by p-value would require a permutation design with sufficient non-ceiling separation, extended altitude data beyond the current horizon, or a larger tail population under altitude controls. None of these are available on the current horizon.

The local $h_1 \times$ sector interaction question is also bounded by the available tail population. P78 rules out the tested binning families and tail thresholds, not every conceivable higher-order relation. A larger tail population could provide enough cell-level resolution to test weaker modulation forms that the current instrument cannot detect.

\subsection{Observed-Horizon Scope}

All continuation claims are confined to the observed terminal band. The row-sector structure and the $h_1$ orthogonality result are empirical laws on this observed horizon. No asymptotic extension is claimed.

\section{Limitations}

The primary claims are empirical and confined to the observed horizon. The sector-amplitude leadership of modulus $1009$ and co-significance of modulus $990$ are properties of the terminal band tested here; behavior at greater heights would require fresh data.

The permutation ceiling affects p-value interpretation for the strongest sector configurations. P-value ranking between modulus $990$ and modulus $1009$ cannot be used to crown a unique winner on the current evidence.

The local $h_1$ interaction tests operate on sparse cells. With the current tail population divided across base-6 sectors and local $h_1$ bins, weak interactions may be undetectable. P78 rules out strong bin-resolution interaction in the tested families; it does not rule out weak local phase conditioning below the current instrument resolution.

The $\sim 10^3$-row sector seam was identified empirically through the P74 window-scale sweep. The tested moduli are reported as part of the audit protocol, and claims are phrased accordingly: modulus $1009$ is an amplitude leader among tested decompositions, not a uniquely proven physical modulus.

\section{Conclusion}

Terminal adjacent-zero spacing residuals exhibit a row-sector rebound structure, amplitude-dominant under modulus $1009$, co-significant under modulus $990$, and empirically orthogonal to the previously published $h_1$ packet organization on the observed terminal horizon. Five sequential audits establish this structure and bound what it is not: the rebound is not gated by $h_1$ corridor membership, local $h_1 \times$ sector modulation is not confirmed, and the apparent growth of interaction strength under finer binning is explained by sparsity inflation. The continuation claim is deliberately narrow: two empirical structures coexist in the terminal band on the observed horizon, and the tested coupling forms between them do not survive null controls.

\section*{Data and Artifact Availability}

All audit outputs, manifests, the frozen parameter registry, the literal claim evidence matrix, and the hardened safe-language document are archived in \texttt{ForgeV16c\_P79B\_LEDGER\_HYGIENE\_outputs} on the author's Drive, with timestamps as recorded in the P79B manifest (\texttt{\p79btimestamp}). The master parquet and seam-stitched board are as described in the principal paper.

\section*{Use of AI Tools}

Generative AI tools were used for audit design review, code critique, adversarial claim evaluation, and draft assistance. All research design decisions, audit execution, interpretation, and final manuscript responsibility remain with Jeffery Huckstead.

\section*{Author Contributions}

Jeffery Huckstead conceived and executed the audit chain, designed the claim ledger and freeze protocol, and wrote the manuscript.

\section*{Acknowledgments}

The author thanks the structured audit-review and adversarial critique process that shaped the ForgeV16c analytical chain.

\begin{thebibliography}{9}

\bibitem{huckstead_h1_principal}
J. Huckstead,
\emph{Phase-Locked Extreme Compression in Adjacent Zeta-Zero Spacing: A Seam-Stitched Empirical Study Through the Full Available LMFDB Horizon},
Zenodo preprint, 2026.
DOI: \href{https://doi.org/\principaldoi}{\principaldoi}.

\bibitem{riemann1859}
B. Riemann,
\emph{Ueber die Anzahl der Primzahlen unter einer gegebenen Grösse},
Monatsberichte der Berliner Akademie, 1859.

\bibitem{odlyzko_tables}
A. M. Odlyzko,
\emph{Tables of zeros of the Riemann zeta function},
available at \url{http://www.dtc.umn.edu/~odlyzko/zeta_tables/}, accessed 2026-05-03.

\bibitem{montgomery1973}
H. L. Montgomery,
\emph{The pair correlation of zeros of the zeta function},
Analytic Number Theory, Proceedings of Symposia in Pure Mathematics, vol. 24, American Mathematical Society, 1973.

\bibitem{mehta_random_matrices}
M. L. Mehta,
\emph{Random Matrices},
3rd ed., Elsevier, 2004.

\end{thebibliography}

\end{document}
